
Current alignment evaluation measures whether AI outputs match human preferences without considering whether excessive alignment affects human oversight through automation bias. This work presents a measurement framework for bidirectional alignment dynamics as coupled variables (ACE and HOR) with experimental protocols for testing coupling mechanisms via capability-invariant policy-layer manipulation.

Our infrastructure validation demonstrates methodological rigor: mock data detection prevented false validation, real dataset verification ensures reproducibility, and dependency-driven hypothesis decomposition with gate-based validation reduces researcher degrees of freedom. However, no empirical evidence exists. All predictions and assumptions remain inconclusive/unverified pending experiment execution (~$1,620 API cost, ~4 hours runtime for h-e1). Our contribution is methodological infrastructure, not empirical findings.

The immediate next step is h-e1 execution to test capability invariance. If h-e1 passes, proceed to mechanism studies (h-m1 through h-m4) requiring human participants and longitudinal data collection. If h-e1 fails, redesign ACE operationalization. Beyond empirical validation, extend the framework to alternative alignment architectures and multi-agent collaboration.

Measurement infrastructure is a necessary foundation for empirical research in emerging problem spaces. We provide the tools for measuring bidirectional alignment dynamics. Future work executes the experiments.
